% Source: https://artofproblemsolving.com/community/c3h1574565_2009_mathcounts_nationals_number_16
% Stable identity: aops-topic-1574565-q01
Let there be $J$ juniors and $S$ seniors. The condition says
\[
\dfrac{3}{5}J=\dfrac{2}{3}\cdot\dfrac{6}{7}S=\dfrac{4}{7}S,\qquad\text{so}\qquad J=\dfrac{5}{3}\cdot\dfrac{4}{7}S=\dfrac{20}{21}S .
\]
The number who passed is $\dfrac{4}{7}S+\dfrac{6}{7}S=\dfrac{10}{7}S$, and the group has $J+S=\dfrac{41}{21}S$ students. The fraction who passed is
\[
\dfrac{\dfrac{10}{7}S}{\dfrac{41}{21}S}=\dfrac{10}{7}\cdot\dfrac{21}{41}=\boxed{\dfrac{30}{41}}.
\]
(For example, $21$ seniors and $20$ juniors: $18$ seniors and $12$ juniors pass, $30$ of $41$.)
